% =============================================================
% 子文件模板：题目描述部分
% 用法：在主文件中使用 \input{filename.tex} 导入
% =============================================================

\begin{center}
    % 建议格式：序号. 中文题目名（英文题目名）
    \section{C. 哥的拔河猜想 (theory)}
\end{center}

\subsection{题目描述}

% 在此处填写题目背景和详细描述
% 行内公式请使用 $...$，例如 $a+b$

**请注意：本题与真正的哥德巴赫猜想没有直接关系，哥只是想拔河。**

在数学界，有一个著名的猜想叫做**哥德巴赫猜想**。它大意是：任意一个大于 $2$ 的偶数，都可以表示成两个质数之和。

例如：

$$
10=3+7
$$

$$
18=5+13
$$

不过，哥觉得两个质数拔河未免太孤单了。

于是，哥提出了自己的猜想：如果让连续若干个质数站成一队，它们的和也许能拔出一些有趣的数字。

我们称两个质数是**相邻的质数**，当且仅当它们之间不存在其他质数。

例如，$2$ 和 $3$ 是相邻的质数，$3$ 和 $5$ 是相邻的质数，$7$ 和 $11$ 也是相邻的质数。

现在哥给了你两个正整数 $N,k$，你需要在所有不超过 $N$ 的正整数中，找出能够表示成**连续 $k$ 个质数之和**的最大数。

如果不存在这样的数，则输出 `-1`。

更形式化地，设质数序列为：

$$
p_1=2,p_2=3,p_3=5,\dots
$$

你需要找到最大的整数 $S$，满足 $S\le N$并且存在正整数 $i$，使得：

$$
S=p_i+p_{i+1}+p_{i+2}+\cdots+p_{i+k-1}
$$

如果不存在满足条件的 $S$，则输出 `-1`。

\subsection{输入格式}

% 描述输入数据的格式，例如：第一行包含一个整数 $n$。
第一行包含一个正整数 $T$，表示数据组数。

接下来 $T$ 行，每行包含两个正整数 $N,k$，表示一组询问。

\subsection{输出格式}

% 描述输出数据的格式
对于每组询问，输出一行一个整数，表示不超过 $N$ 且能够表示成连续 $k$ 个质数之和的最大数。

如果不存在这样的数，则输出 `-1`。

\subsection{样例输入1}

% 将样例输入直接粘贴在下方，无需额外格式
\begin{lstlisting}
3
20 2
20 3
20 4
\end{lstlisting}

\subsection{样例输出1}

% 将样例输出直接粘贴在下方
\begin{lstlisting}
18
15
17
\end{lstlisting}

\subsection{样例解释1}

对于第一组询问，$N=20,k=2$。

连续 $2$ 个质数之和依次为：

\begin{lstlisting}
2+3=5
3+5=8
5+7=12
7+11=18
11+13=2
\end{lstlisting}

其中不超过 $20$ 的最大值为 $18$。

对于第二组询问，$N=20,k=3$。

连续 $3$ 个质数之和依次为：

\begin{lstlisting}
2+3+5=10
3+5+7=15
5+7+11=23
\end{lstlisting}

其中不超过 $20$ 的最大值为 $15$。

对于第三组询问，$N=20,k=4$。

连续 $4$ 个质数之和依次为：

\begin{lstlisting}
2+3+5+7=17
3+5+7+11=26
\end{lstlisting}

其中不超过 $20$ 的最大值为 $17$。


\subsection{数据范围和约定}

% 描述数据规模和特殊限制
% 示例：对于 $100\%$ 的数据，$1 \leq n \leq 100$。
对于 $20\%$ 的数据，$1 \le N \le 100$。

对于另外 $20\%$ 的数据，$T = 1$。

对于另外 $20\%$ 的数据，所有询问中的 $N$ 均相等。

对于 $100\%$ 的数据，$1 \le T < 2000$，$1 \le N \le 10^6$，$1 \le k$。

